\documentclass[aspectratio=169,11pt]{beamer}
\usepackage[T1]{fontenc}
\usepackage[utf8]{inputenc}
\IfFileExists{beamerthemeifsclean.sty}{}{%
 \IfFileExists{brazilian.ldf}{\usepackage[brazilian]{babel}}{\usepackage[english]{babel}}
}
\usepackage{lmodern}
\usepackage{amsmath,amssymb,booktabs,array,graphicx,tikz,ragged2e,url}
\usetikzlibrary{positioning,arrows.meta,calc,fit,backgrounds}
\definecolor{ifscgreen}{RGB}{0,128,0}
\definecolor{darkgreen}{RGB}{0,90,0}
\definecolor{lightgraybox}{RGB}{245,245,245}
\IfFileExists{beamerthemeifsclean.sty}{\usetheme{ifsclean}}{%
 \usetheme{Boadilla}
 \usecolortheme[named=darkgreen]{structure}
 \setbeamercolor{title}{fg=white,bg=darkgreen}
 \setbeamercolor{frametitle}{fg=darkgreen,bg=white}
 \setbeamercolor{block title}{fg=white,bg=darkgreen}
 \setbeamercolor{block body}{bg=lightgraybox}
 \setbeamertemplate{navigation symbols}{}
 \setbeamertemplate{footline}{\leavevmode\hbox{%
 \begin{beamercolorbox}[wd=.76\paperwidth,ht=2.8ex,dp=1.2ex,leftskip=1em]{author in head/foot}\usebeamerfont{author in head/foot}Wagner Flores dos Santos\quad TCC1\end{beamercolorbox}%
 \begin{beamercolorbox}[wd=.24\paperwidth,ht=2.8ex,dp=1.2ex,rightskip=1em plus1fil]{date in head/foot}\hfill 2026.2\quad\insertframenumber\end{beamercolorbox}}}
}
\setbeamersize{text margin left=0.8cm,text margin right=0.8cm}
\setlength{\emergencystretch}{2em}
\newcommand{\fonte}[1]{\par\vspace{0.5em}{\tiny\color{gray}\textbf{Fonte:} #1}}
\newcommand{\refweb}[2]{\href{#1}{\textcolor{darkgreen}{#2}}}
\tikzset{caixa/.style={draw=darkgreen,rounded corners=2pt,fill=green!5,align=center,minimum height=1cm,text width=2.4cm,font=\small},seta/.style={-{Latex[length=2mm]},thick,darkgreen}}
\title[Agente único e multiagente]{Análise Comparativa entre Arquiteturas de Agente Único e \\ Multiagente Baseadas em LLM para Chatbots de Atendimento}
\subtitle{Seminário de apresentação do tema -- TCC1}
\author{Wagner Flores dos Santos}
\institute{Engenharia de Telecomunicações -- IFSC Câmpus São José\\Orientador: Prof. Arliones Stevert Hoeller Júnior, Dr.}
\date{2026.2}
\newcommand{\origem}{\fonte{Pré-projeto do autor, versão disponibilizada em 14/11/2026.}}
\begin{document}
% Capa: 20 s. Conteúdo principal: 8 min 40 s. Total planejado: 9 min.
\begin{frame}[plain,noframenumbering]\maketitle\end{frame}
%% Tempo: 55 s
\begin{frame}{Especializar agentes pode melhorar o atendimento, mas tem custo}
\begin{block}{Problema de pesquisa}
Como agente único e multiagente diferem em qualidade, encaminhamentos e consumo de recursos ao longo de uma conversa?
\end{block}
\vspace{.25cm}
\begin{itemize}
\item LLM: modelo de linguagem em grande escala.
\item O cliente informa dados aos poucos e pode mudar de assunto.
\item Mais agentes acrescentam decisões e comunicação ao atendimento.
\end{itemize}
\origem
\end{frame}

\begin{frame}{A escolha da arquitetura afeta qualidade e recursos}
\begin{columns}[T]
\column{.48\textwidth}
\begin{block}{Agente único}
Reúne as regras das quatro áreas.\\[.3cm]
Menos coordenação interna; contexto mais amplo.
\end{block}
\column{.48\textwidth}
\begin{block}{Multiagente}
Distribui regras entre especialistas.\\[.3cm]
Especialização; novas decisões e possível perda de contexto.
\end{block}
\end{columns}
\vspace{.35cm}
\textbf{Justificativa:} medir se os ganhos observados compensam chamadas e tokens adicionais.
\fonte{Síntese do pré-projeto; Yan et al. (2026); Amri et al. (2026).}
\end{frame}

\begin{frame}{Os estudos revisados deixam a avaliação conjunta em aberto}
\small
\renewcommand{\arraystretch}{1.35}
\begin{tabular}{@{}p{3cm}p{4.25cm}p{6cm}@{}}\toprule
Estudo & Contribuição & Limite para esta proposta\\\midrule
Moncada et al. (2025) & Classificação e respostas a chamados & Sem comparação com agente único ou conversa com vários turnos\\
Lee et al. (2025) & Arquitetura em telecomunicações & Sem comparação quantitativa própria\\
Amri et al. (2026) & Compara arquiteturas em diálogos & Foco na transferência para humanos; sem consumo de recursos\\\bottomrule
\end{tabular}
\vspace{.3cm}
\begin{block}{Lacuna no conjunto revisado}
Avaliar conjuntamente qualidade, retenção, conclusão, encaminhamentos e consumo em conversas com vários turnos.
\end{block}
\fonte{Síntese da revisão do pré-projeto, seções 2.2 e 2.3.}
\end{frame}

\begin{frame}{O objetivo é comparar as arquiteturas sob condições comuns}
\begin{block}{Objetivo geral}
Comparar experimentalmente agente único e multiagente no atendimento textual de um provedor de Internet.
\end{block}
\vspace{.2cm}
\begin{enumerate}
\item Definir base e regras comuns; implementar dois protótipos.
\item Elaborar e revisar os cenários conversacionais.
\item Avaliar respostas, conclusão e encaminhamentos.
\item Medir chamadas, tokens e custo; comparar as diferenças.
\end{enumerate}
\origem
\end{frame}

\begin{frame}{Um único agente concentra as regras das quatro áreas}
\centering\resizebox{.96\textwidth}{!}{%
\begin{tikzpicture}[
        font=\small,
        bloco/.style={
            draw,
            rounded corners=3pt,
            align=center,
            fill=blue!6,
            minimum height=11mm,
            text width=24mm
        },
        auxiliar/.style={
            bloco,
            fill=gray!8,
            text width=22mm
        },
        agente/.style={
            bloco,
            fill=blue!12,
            text width=21mm,
            minimum height=16mm
        },
        troca/.style={
            <->,
            >=Latex,
            thick
        },
        registro/.style={
            ->,
            >=Latex,
            thick
        },
        base/.style={
            ->,
            >=Latex,
            dashed,
            gray!80
        }
    ]
    % Robô simplificado.
    \newcommand{\robo}{%
        \begin{tikzpicture}[
            x=1mm,
            y=1mm,
            baseline=-0.5mm,
            line width=0.5pt
        ]
            \draw[fill=blue!15,rounded corners=1pt]
                (-3,-2) rectangle (3,2);
            \fill (-1.2,0.4) circle (0.35);
            \fill (1.2,0.4) circle (0.35);
            \draw (-1,-0.9) -- (1,-0.9);
            \draw (0,2) -- (0,3);
            \fill (0,3) circle (0.35);
            \draw (-3,-0.7) -- (-3.8,-0.7)
                  -- (-3.8,0.7) -- (-3,0.7);
            \draw (3,-0.7) -- (3.8,-0.7)
                  -- (3.8,0.7) -- (3,0.7);
        \end{tikzpicture}%
    }
    \node[bloco] (u1) at (-5,-1.2)
        {Usuário};
    \node[bloco] (o1) at (-1.5,-1.2)
        {Orquestrador\\[-1pt]
         \scriptsize Controle em Python};
    \node[agente] (a1) at (2,-1.2)
        {\robo\\Agente único};
    \node[auxiliar] (b1) at (5.2,-1.2)
        {Base completa\\[-1pt]
         \scriptsize Quatro áreas};
    \node[auxiliar] (h1) at (-3.2,-3.5)
        {Histórico e fatos};
    \node[auxiliar] (r1) at (0.2,-3.5)
        {Registro da execução};
    % Comunicação principal.
    \draw[troca] (u1) -- (o1);
    \draw[troca] (o1) -- (a1);
    % Informações da base.
    \draw[base] (b1) -- (a1);
    % Histórico e registro.
    \draw[troca]
        (o1.south) -- ++(0,-0.45)
        -| (h1.north);
    \draw[registro]
        (o1.south) -- ++(0,-0.45)
        -| (r1.north);
    \end{tikzpicture}}
\par\vspace{.25cm}\small O orquestrador mantém o histórico e os registros; o agente decide e responde.\\[.2cm] Linhas tracejadas: informações disponibilizadas ao agente.
\fonte{TikZ da Figura 1 da monografia do autor, adaptado apenas para apresentação em dois slides.}
\end{frame}

\begin{frame}{A triagem e os especialistas são coordenados pelo orquestrador}
\centering\resizebox{.75\textwidth}{!}{%
\begin{tikzpicture}[
        font=\small,
        bloco/.style={
            draw,
            rounded corners=3pt,
            align=center,
            fill=blue!6,
            minimum height=11mm,
            text width=24mm
        },
        auxiliar/.style={
            bloco,
            fill=gray!8,
            text width=22mm
        },
        agente/.style={
            bloco,
            fill=blue!12,
            text width=21mm,
            minimum height=16mm
        },
        troca/.style={
            <->,
            >=Latex,
            thick
        },
        registro/.style={
            ->,
            >=Latex,
            thick
        },
        base/.style={
            ->,
            >=Latex,
            dashed,
            gray!80
        }
    ]

    \newcommand{\robo}{%
        \begin{tikzpicture}[
            x=1mm,
            y=1mm,
            baseline=-0.5mm,
            line width=0.5pt
        ]
            \draw[fill=blue!15,rounded corners=1pt]
                (-3,-2) rectangle (3,2);
            \fill (-1.2,0.4) circle (0.35);
            \fill (1.2,0.4) circle (0.35);
            \draw (-1,-0.9) -- (1,-0.9);
            \draw (0,2) -- (0,3);
            \fill (0,3) circle (0.35);
            \draw (-3,-0.7) -- (-3.8,-0.7)
                  -- (-3.8,0.7) -- (-3,0.7);
            \draw (3,-0.7) -- (3.8,-0.7)
                  -- (3.8,0.7) -- (3,0.7);
        \end{tikzpicture}%
    }

    \begin{scope}[yshift=-10mm]
        \node[bloco] (u2) at (-5,-5.4)
            {Usuário};
        \node[bloco] (o2) at (-1.5,-5.4)
            {Orquestrador\\[-1pt]
             \scriptsize Controle em Python};
        \node[auxiliar] (h2) at (2,-5.4)
            {Histórico e fatos};
        \node[auxiliar] (r2) at (5.2,-5.4)
            {Registro da execução};
        % Comunicação com usuário e histórico.
        \draw[troca] (u2) -- (o2);
        \draw[troca] (o2) -- (h2);
        % Registro ligado ao orquestrador por cima dos blocos.
        \draw[registro]
            (o2.north) -- ++(0,0.45)
            -| (r2.north);
        % Agentes.
        \node[agente] (t) at (-5.1,-8.2)
            {\robo\\Triagem};
        \node[agente] (p) at (-2.2,-8.2)
            {\robo\\Primeiro\\atendimento};
        \node[agente] (c) at (0.7,-8.2)
            {\robo\\Comercial};
        \node[agente] (v) at (3.6,-8.2)
            {\robo\\Ouvidoria};
        \node[agente] (s) at (6.5,-8.2)
            {\robo\\Suporte\\técnico};
        % Barramento de comunicação do orquestrador.
        \coordinate (centro) at (-1.5,-6.8);
        \draw[thick] (o2.south) -- (centro);
        \draw[thick]
            (-5.1,-6.8) -- (6.5,-6.8);
        \foreach \n in {t,p,c,v,s}{
            \draw[troca]
                (\n.north |- centro) -- (\n.north);
        }
        % Base de conhecimento.
        \node[
            auxiliar,
            text width=55mm
        ] (b2) at (0.7,-10.7)
            {Base comum e bases por área};
        \coordinate (bases) at (0.7,-9.7);
        \draw[dashed,gray!80]
            (b2.north) -- (bases);
        \draw[dashed,gray!80]
            (-5.1,-9.7) -- (6.5,-9.7);
        \foreach \n in {t,p,c,v,s}{
            \draw[base]
                (\n.south |- bases) -- (\n.south);
        }
    \end{scope}
    \end{tikzpicture}}
\par\vspace{.25cm}\scriptsize Conexões: possibilidades de comunicação, sem execução simultânea obrigatória.\\O especialista permanece ativo; mudanças de assunto solicitam nova triagem.\\Linhas tracejadas: base comum e informações por responsabilidade.
\fonte{TikZ da Figura 1 da monografia do autor, adaptado apenas para apresentação em dois slides.}
\end{frame}

\begin{frame}{48 cenários testarão continuidade e mudanças de assunto}
\begin{columns}[T]
\column{.49\textwidth}
\begin{block}{Conjunto de teste}
4 áreas $\times$ 3 situações $\times$ 4 variações de redação = \textbf{48 cenários}.
\end{block}
Cada cenário nas duas arquiteturas:\\[.15cm]
\textbf{96 conversas completas}, uma execução por combinação.
\column{.47\textwidth}
\small
\begin{tabular}{@{}p{4cm}p{2.1cm}@{}}\toprule
Mensagem (exemplo) & Área\\\midrule
Quero conhecer os planos. & Comercial\\
Meu CEP é 88110-000. & Comercial\\
Já sou cliente e estou sem Internet. & Suporte\\
A luz LOS está vermelha. & Suporte\\
Isso se repete; quero reclamar. & Ouvidoria\\\bottomrule
\end{tabular}
\end{columns}
\vspace{.35cm}
\textbf{Dados fictícios.} As variações preservam fatos e critérios de conclusão.
\origem
\end{frame}

\begin{frame}{A avaliação combinará métricas e verificações objetivas}
\renewcommand{\arraystretch}{1.35}
\small
\begin{tabular}{@{}p{3.2cm}p{6.5cm}p{3.5cm}@{}}\toprule
Dimensão & O que observar & Como verificar\\\midrule
Qualidade & Relevância, retenção e completude & DeepEval + revisão de amostra\\
Correção e tarefa & Respeito às regras e condições de conclusão & Critérios por cenário\\
Encaminhamentos & Área por turno e indicação de humano & Acurácia e matriz de confusão; regras\\
Recursos & Chamadas, tokens de entrada e saída; custo & Registros em Python\\\bottomrule
\end{tabular}
\vspace{.35cm}
\textbf{Comparação equivalente:} área identificada no agente único versus área do especialista acionado no multiagente.
\vspace{.2cm}
Trocas internas serão contadas apenas no multiagente; indicação de humano, nas duas arquiteturas.
\fonte{Pré-projeto, seção 3.7; Confident AI (2026).}
\end{frame}
%% Tempo: 65 s
\begin{frame}{O piloto fixará as condições antes da comparação}
\begin{columns}[T]
\column{.53\textwidth}
\begin{block}{Controles previstos}
Mesmo GPT-4.1-mini, parâmetros, conteúdo total, regras e mensagens.\\[.2cm]
Mesmo avaliador e limiares nas duas arquiteturas.
\end{block}
\textbf{Sequência:} revisar $\rightarrow$ piloto $\rightarrow$ congelar versões $\rightarrow$ executar $\rightarrow$ avaliar.
\column{.43\textwidth}
\begin{block}{Limites da interpretação}
Uma execução não mede a variabilidade da mesma entrada.\\[.2cm]
48 versões derivam de 12 situações-base.\\[.2cm]
Resultados restritos ao modelo e aos cenários escolhidos.
\end{block}
\end{columns}
\vspace{.25cm}
Python, arquivos JSON e Git para implementar, registrar e versionar. O piloto fica fora dos resultados finais.
\origem
\end{frame}

\begin{frame}{O TCC2 entregará evidências para a escolha da arquitetura}
\small
\renewcommand{\arraystretch}{1.25}
\begin{tabular}{@{}p{7cm}cccccc@{}}\toprule
Atividade em 2027 & Fev & Mar & Abr & Mai & Jun & Jul\\\midrule
Base, prompts, protótipos e cenários & X & X & & & &\\
Métricas e piloto & & X & X & & &\\
Congelamento e experimento & & & X & & &\\
Análise comparativa & & & X & X & &\\
Redação e revisões & X & X & X & X & X &\\
Preparação e defesa & & & & & X &\\
Correções e entrega & & & & & X & X\\\bottomrule
\end{tabular}
\begin{block}{Mensagem central}
A especialização pode melhorar o atendimento, mas acrescenta coordenação. O estudo medirá qualidade e recursos para identificar as diferenças, sem pressupor um vencedor.
\end{block}
\fonte{Cronograma previsto no pré-projeto; defesa e entrega sujeitas ao calendário oficial.}
\end{frame}
\appendix

\begin{frame}{Referências utilizadas na apresentação}
\scriptsize
\begin{itemize}
\item AMRI, M. F. et al. Optimizing LLM Chatbot to Human Handover with Multi-Agent AI Approach. ICADEIS, 2026, p. 1--6. DOI: 10.1109/ICADEIS71120.2026.11644567.
\item CONFIDENT AI. DeepEval Documentation. 2026. \url{https://deepeval.com/docs/introduction}. Acesso: 8 out. 2026.
\item LEE, K.-Y.; TSENG, P.-A.; CHEN, W.-C. Enhancing Telecom Operation Support Systems with Multi-Agent Large Language Models. ICCAI, 2025, p. 309--316. DOI: 10.1109/ICCAI66501.2025.00055.
\item MONCADA, C. et al. Large Language Models and Multi-Agent Systems for Customer Support Ticket Automation. ETCM, 2025, p. 1--6. DOI: 10.1109/ETCM67548.2025.11304425.
\item YAN, B. et al. Beyond Self-Talk: A Communication-Centric Survey of LLM-Based Multi-Agent Systems. Frontiers of Computer Science, 2026. DOI: 10.1007/s11704-026-50857-y.
\item YANG, W.; XU, P.; ZHANG, X. A Two-Stage Intent Recognition Framework in Multi-Agent Environments. ICCEA, 2025, p. 1--5. DOI: 10.1109/ICCEA65460.2025.11103313.
\item OWOTOGBE, J. Assessing and Enhancing the Robustness of LLM-Based Multi-Agent Systems Through Chaos Engineering. CAIN, 2025, p. 250--252. DOI: 10.1109/CAIN66642.2025.00039.
\item OPENAI. GPT-4.1 mini Model. 2026. \url{https://developers.openai.com/api/docs/models/gpt-4.1-mini}. Acesso: 10 out. 2026.
\end{itemize}
\fonte{Dados bibliográficos reproduzidos da monografia fornecida.}
\end{frame}

\begin{frame}{Reserva: escopo e decisões pendentes}
\begin{itemize}
\item Atendimento textual, quatro áreas, dados sintéticos e documentos no contexto; sem recuperação dinâmica de documentos.
\item Sem clientes reais, voz, multimodalidade, ajuste fino, comparação entre modelos ou injeção de falhas.
\item Latência não será avaliada; a transferência humana será apenas registrada.
\item Versão exata do modelo, parâmetros, avaliador e limiares: registrar e congelar após o piloto.
\item TCC1: fundamentação e planejamento; experimento definitivo no TCC2.
\end{itemize}
\origem
\end{frame}

\begin{frame}{Reserva: custo e critérios de sucesso}
\[C=T_e P_e + T_s P_s\]
\begin{itemize}
\item $T_e,T_s$: tokens de entrada e saída; $P_e,P_s$: preço por token.
\item Registrar fonte e data dos preços no experimento.
\item Concluir uma tarefa pode significar indicar humano com os dados necessários, quando isso for previsto no cenário.
\item Sucesso da pesquisa: comparação executada e diferenças analisadas, independentemente da arquitetura com melhor desempenho.
\end{itemize}
\origem
\end{frame}

\begin{frame}{Reserva: os seis estudos fundamentam o recorte da pesquisa}
\footnotesize
\renewcommand{\arraystretch}{1.15}
\begin{tabular}{@{}p{2.9cm}p{5cm}p{5.3cm}@{}}
\toprule
Estudo & Contribuição & Limite para esta comparação\\
\midrule
Moncada et al. (2025) & Classificação e geração de respostas a chamados & 20 chamados; sem agente único ou conversa com vários turnos\\
Lee et al. (2025) & Coordenação e especialistas em telecomunicações & Sem avaliação quantitativa própria entre arquiteturas\\
Yang et al. (2025) & Seleção de agentes por TF-IDF e LLM & 30 consultas; sem avaliação da conversa após a seleção\\
Amri et al. (2026) & Compara agente único e multiagente em diálogos & Foco na transferência humana; sem consumo de recursos\\
Yan et al. (2026) & Revisão centrada na comunicação entre agentes & Sem experimento próprio comparando as arquiteturas\\
Owotogbe (2025) & Proposta para investigar robustez e propagação de falhas & Validação em andamento; sem resultados quantitativos do framework\\
\bottomrule
\end{tabular}
\par\vspace{.2cm}
\textbf{Recorte do TCC:} qualidade, retenção, conclusão, encaminhamentos e recursos avaliados conjuntamente.
\fonte{Síntese da revisão da monografia, seções 2.2 e 2.3. Referências completas no slide anterior.}
\end{frame}

\begin{frame}{Reserva: DeepEval avalia; Synthesizer apoia os cenários}
\small
\begin{columns}[T]
\column{.48\textwidth}
\begin{block}{Antes dos testes: Synthesizer}
\begin{itemize}
\item Usa contextos ou documentos para sugerir cenários e desfechos esperados.
\item O pesquisador revisa as sugestões, prepara os turnos e as variações de redação.
\item As mensagens são fixadas e aplicadas igualmente às duas arquiteturas.
\end{itemize}
\end{block}
\column{.48\textwidth}
\begin{block}{Após os testes: DeepEval}
\begin{itemize}
\item Relevância: a resposta atende ao assunto do turno?
\item Retenção: preserva os fatos informados na conversa?
\item Completude: atende às necessidades apresentadas?
\end{itemize}
\end{block}
\end{columns}
\vspace{.15cm}
\footnotesize
\textbf{Controles:} revisão humana dos cenários; mesmo avaliador e limiares; revisão manual de uma amostra das avaliações. Encaminhamentos, chamadas e tokens: verificações em Python.
\par\vspace{.15cm}
\textbf{Distinção:} gerar cenários não é validá-los. Neste estudo, o Synthesizer não decidirá sozinho o roteiro final nem substituirá a revisão do pesquisador.
\fonte{Pré-projeto, seções 3.4 e 3.7; documentação DeepEval: \href{https://deepeval.com/docs/golden-synthesizer}{Synthesizer} e \href{https://deepeval.com/docs/metrics-introduction}{métricas}. Consulta: 14/11/2026.}
\end{frame}
\end{document}
